\section*{Historical Note}
\phantomsection
\addcontentsline{toc}{section}{Historical Note}

(N.B. The roman numerals refer to the Bibliography at the end of this note.)

As we have seen (Historical Note to Chapters I-II-III), the problems that led to the integration of differential equations were among the first to be treated by the founders of the infinitesimal calculus of the XVII \(^{th}\) century (notably Descartes and Barrow). Since then the theory of differential equations has not ceased to tax the sagacity of mathematicians, and to be a favoured terrain for the application of the most varied methods of Analysis; the questions that it provokes are very far from being all solved, and the interest that attaches is accordingly so sustained that it constitutes one of the most permanent and fruitful points of contact between mathematics and the experimental sciences: these latter often find a valuable help there, and in exchange constantly provide new problems.

Of the numerous chapters which a modern complete study of differential equations should encompass, we have not wished to expound here but the two most elementary, treating existence theorems and linear equations, the variable being assumed to be real. This is why we also limit our brief historical account to these two aspects. From the beginning of the XVIII \(^{th}\) century mathematicians were convinced that the ``general'' integral of a differential equation of order n depends on n arbitrary constants, and that ``in general'' there exists one and only one integral with given values for the function and its first n - 1 derivatives at a given value \(x_{0}\) of the variable: a conviction which they justified by the process (going back to Newton) of calculating one by one the coefficients of the Taylor expansion of the solution at the point \(x_{0}\) , with the help of the differential equation itself and the first n coefficients. But up until Cauchy no one had studied the convergence of the series so obtained, nor proved that its sum was a solution of the differential equation; and of course, there was no question of nonanalytic differential equations. Among the various methods invented by Cauchy to prove the existence of integrals of differential equations, that which we have followed ((IV) and (IV bis)), generalized a little later by Lipschitz, is particularly interesting for the case of nonanalytic equations and approximation of integrals.

The linear differential equations were among the first to draw attention. Leibniz and Jakob Bernoulli integrated the first order linear equation by two quadratures ((I), v. II, p.~731); the general integral of the arbitrary linear equation with constant coefficients and arbitrary right hand side was obtained by Euler ((II), p.~296-354); d'Alembert solved the linear systems with constant coefficients in the same way. A little later, Lagrange (III) treated the general theory of linear equations of order \(n\) , recognised that the integral of the homogeneous equation is a linear function of the \(n\) constants of integration, introduced the adjoint equation, discovered the reduction of the order of the homogeneous equation when one knows particular solutions, and the method of variation of parameters to solve the inhomogeneous equation. The obscure points in Lagrange's theory (notably in what concerns the linear independence of the integrals) were clarified by Cauchy, whose exposition (IV bis) has remained quasi definitive, apart from the improvements in detail provided by matrix notation and the theory of elementary divisors.

